\section{Casos sintéticos: planteamiento y soluciones de referencia}
\label{app:casos-sinteticos}

Este apéndice especifica, para cada caso de prueba con datos sintéticos
(Exp.~1 a~5), el planteamiento matemático del problema --- ecuaciones de
gobierno, dominio, condiciones de borde e iniciales y batimetría --- y su
solución de referencia. Para los casos con solución analítica cerrada
(Exp.~1, 2 y 5) se da la expresión exacta; para los casos sin forma cerrada
(Exp.~3 y 4) la referencia es la solución numérica de un esquema de volúmenes
finitos. Las ecuaciones de aguas someras (SWE) y su notación se introdujeron
en la Sección~\ref{sec:problema}; aquí se particularizan a cada caso. La
gravedad es $g = 9.81$\,m\,s$^{-2}$.

\subsection{Exp.\ 1 --- Bump subcrítico 1D}
\label{app:exp1}

\paragraph{Problema.}
Flujo 1D estacionario sin fricción sobre un bump parabólico, gobernado por las
SWE estacionarias (caudal específico $q = h u$ constante y conservación de la
energía). La batimetría es
\begin{equation}
  \zb(x) =
  \begin{cases}
    H_b - \dfrac{H_b}{w^2}\,(x - x_0)^2, & |x - x_0| < w,\\[1.2ex]
    0, & \text{en otro caso,}
  \end{cases}
\end{equation}
con altura $H_b = 0.2$\,m, semiancho $w = 2$\,m y centro $x_0 = 0$, en el
dominio $[-10, 10]$\,m. Las condiciones de borde son el caudal específico
$q = 4.42$\,m$^2$/s aguas arriba y la profundidad $h_{\text{down}} = 2.0$\,m
aguas abajo.

\paragraph{Solución de referencia (analítica).}
En flujo estacionario sin fricción la energía específica se conserva a lo
largo del canal (ecuación de Bernoulli):
\begin{equation}
  \frac{q^2}{2 g\,h(x)^2} + h(x) + \zb(x) = C,
  \label{eq:app-bernoulli}
\end{equation}
donde la constante $C$ se fija con la condición aguas abajo,
$C = q^2/(2 g\,h_{\text{down}}^2) + h_{\text{down}} + \zb(x_{\text{fin}})$.
Multiplicando \eqref{eq:app-bernoulli} por $h^2$ se obtiene, en cada $x$, la
ecuación cúbica
\begin{equation}
  h^3 + \bigl(\zb(x) - C\bigr)\,h^2 + \frac{q^2}{2 g} = 0,
\end{equation}
de cuyas raíces se toma la \emph{subcrítica}: real, positiva y mayor que la
profundidad crítica $h_c = (q^2/g)^{1/3}$. La velocidad y la superficie libre
son $u(x) = q/h(x)$ y $\eta(x) = h(x) + \zb(x)$. El planteamiento sigue a
Delestre et al.\ \cite{delestre2013} (compilación SWASHES) y a Dazzi et al.\
\cite{dazzi2024}.

\subsection{Exp.\ 2 --- Cuenca de Thacker 1D}
\label{app:exp2}

\paragraph{Problema.}
Oscilación no lineal de una masa de agua en una cuenca parabólica cerrada,
gobernada por las SWE 1D transitorias sin fricción. La batimetría es
\begin{equation}
  \zb(x) = h_0\left(\frac{x^2}{a^2} - 1\right),
\end{equation}
con profundidad característica $h_0 = 0.5$\,m y semiancho $a = 1$\,m, en el
dominio $[-2, 2]$\,m. El problema es cerrado: no hay condiciones de borde de
flujo, y la condición inicial es la solución analítica evaluada en $t = 0$.

\paragraph{Solución de referencia (analítica).}
Thacker \cite{thacker1981} obtuvo una solución exacta de superficie libre
plana con línea de costa móvil. La frecuencia angular y el período de
oscilación son
\begin{equation}
  \omega = \frac{\sqrt{2 g h_0}}{a},
  \qquad T = \frac{2\pi}{\omega} \approx 2.006~\text{s},
\end{equation}
y los campos
\begin{align}
  h(x,t)    &= \max\!\left(0,\;
               h_0\left[1 - \left(\frac{x + \tfrac{1}{2}\cos(\omega t)}{a}\right)^{\!2}\,\right]\right),\\
  u(x,t)    &= \tfrac{1}{2}\,\omega\,\sin(\omega t)
               \quad\text{donde } h > 0\ \ (\text{0 en seco}),\\
  \eta(x,t) &= h(x,t) + \zb(x).
\end{align}
La superficie libre permanece plana y oscila en el tiempo, mientras la línea
de costa se desplaza. El caso está también compilado en SWASHES
\cite{delestre2013}.

\subsection{Exp.\ 3 --- Dos cilindros 2D}
\label{app:exp3}

\paragraph{Problema.}
Flujo 2D transitorio sobre dos obstáculos cilíndricos sólidos, gobernado por
las SWE 2D en forma conservativa,
$\bm{U}_t + \bm{F}(\bm{U})_x + \bm{G}(\bm{U})_y = \bm{S}(\bm{U}, \zb)$ con
$\bm{U} = (h, hu, hv)$. La configuración se replica fielmente del benchmark
2D de Ruppenthal y Kuzmin \cite{ruppenthal2026} (Sección~7.2). La batimetría
consiste en dos cilindros con paredes verticales (sin suavizado),
\begin{equation}
  \zb(x,y) =
  \begin{cases}
    0.2\,\text{m}, & \text{si } \sqrt{(x-8)^2 + (y-8)^2} \le 4\,\text{m}, \\
    0.3\,\text{m}, & \text{si } \sqrt{(x-15)^2 + (y-15)^2} \le 2\,\text{m}, \\
    0,             & \text{en otro caso,}
  \end{cases}
\end{equation}
en el dominio $[0, 25]^2$\,m. Los dos cilindros (alturas $H_1 = 0.2$ y
$H_2 = 0.3$\,m, radios $R_1 = 4$ y $R_2 = 2$\,m) están alineados sobre la
diagonal $x = y$, paralelos a la dirección del flujo inicial. La condición
inicial corresponde a superficie libre en reposo y velocidad uniforme:
\begin{equation}
  \eta(x,y,0) = h + \zb = 2\,\text{m}, \qquad
  \bm{v}(x,y,0) = (2.21,\ 2.21)\,\text{m/s},
\end{equation}
de modo que la profundidad inicial $h(x,y,0) = 2 - \zb(x,y)$ es menor sobre
los cilindros que sobre el lecho plano. Tiempo final $T = 60$\,s con paso
$\Delta t = 10^{-2}$\,s.

\paragraph{Solución de referencia (numérica).}
Este caso no admite solución analítica cerrada. La referencia se obtiene
resolviendo las SWE 2D sobre una malla cartesiana de $50\times 50$ celdas
($\Delta x = \Delta y = 0.5$\,m), con un esquema de volúmenes finitos, flujos
numéricos de tipo HLL (Harten--Lax--van~Leer) sobre estados reconstruidos
por la técnica hidrostática de Audusse \cite{audusse2004}, well-balanced
para topografía discontinua; el término fuente topográfico queda absorbido
en correcciones de presión en las interfaces, siguiendo en lo demás la
formulación estándar de LeVeque \cite{leveque2002}. Las condiciones de
frontera son Dirichlet de inflow en las dos caras por las que entra el
flujo --$x = 0$ y $y = 0$-- prescribiendo el estado uniforme
$(h, u, v) = (\eta_0, u_0, v_0)$ vía celdas fantasma; las caras opuestas
($x = L_x$, $y = L_y$) usan extrapolación de gradiente cero (outflow).
Ruppenthal y Kuzmin no especifican BCs explícitamente para §7.2; ésta
es la elección más natural para preservar la corriente uniforme
prescrita en la CI a lo largo de los 60\,s simulados. Ruppenthal y Kuzmin emplean un esquema de
elementos finitos con limiters de flujo (MCL) sobre la misma malla; ambas
son discretizaciones estándar y convergentes de las SWE.

\subsection{Exp.\ 4 --- Topografía variable 2D (Tian dT10)}
\label{app:exp4}

\paragraph{Problema.}
Flujo 2D transitorio sobre una topografía oscilatoria, gobernado por las SWE
2D en forma conservativa. La configuración se replica fielmente del caso
dinámico \emph{dT10} de Tian et al.\ \cite{tian2025} (Sección~3.3),
denominado allí ``problema mareal'' por analogía con la geometría de fondo
(no por la presencia de forzamiento mareal externo). La topografía es
\begin{equation}
  \zb(x,y) = 1 + 0{.}01\,\cos\!\left(\frac{\pi x}{2}\right)
                          \cos\!\left(\frac{\pi y}{2}\right),
\end{equation}
en el dominio $(x,y) \in [-2, 2]^2$\,m. Tiene máximos en el centro y en las
cuatro esquinas ($\zb = 1{.}01$\,m) y mínimos en los puntos medios de los
bordes ($\zb = 0{.}99$\,m). La condición inicial es desbalanceada respecto a
la superficie libre:
\begin{equation}
  h(x,y,0) = \zb(x,y), \qquad u(x,y,0) = v(x,y,0) = 0,
\end{equation}
de modo que la superficie libre inicial $\eta(x,y,0) = h + \zb = 2\zb$ hereda
el mismo patrón de máximos y mínimos. Bajo gravedad el agua fluye desde las
regiones altas (centro y esquinas) hacia las bajas (mitades de los bordes).
Las condiciones de borde son \emph{periódicas} en $x$ y $y$. Tiempo final
$T = 0{.}5$\,s.

\paragraph{Solución de referencia (numérica).}
Sin forma cerrada. Tian resuelve dT10 con el esquema entropy-stable ES1 de
Fjordholm et al.\ (2011) sobre una malla muy fina
($\Delta x = \Delta y = 0{.}01$\,m, CFL $= 0{.}25$). En esta tesis la
referencia se obtiene con el mismo solver FV-HLL del Exp.~3
(LeVeque \cite{leveque2002}) con condiciones de borde periódicas. Ambos
son esquemas estándar y convergentes para SWE con topografía variable.

\subsection{Exp.\ 5 --- Paraboloide de Thacker 2D}
\label{app:exp5}

\paragraph{Problema.}
Oscilación axisimétrica de una masa de agua en un paraboloide de revolución,
gobernada por las SWE 2D transitorias sin fricción. La batimetría es
\begin{equation}
  \zb(x,y) = -h_0\left(1 - \frac{r^2}{a^2}\right),
  \qquad r = \sqrt{(x - x_c)^2 + (y - y_c)^2},
\end{equation}
con $h_0 = 0.1$\,m, $a = 1$\,m y centro $(x_c, y_c) = (2, 2)$\,m, en el dominio
$[0, 4]^2$\,m. El problema es cerrado, con una frontera seco-mojado móvil.

\paragraph{Solución de referencia (analítica).}
Thacker \cite{thacker1981} obtuvo la solución exacta axisimétrica (compilada
en SWASHES \cite{delestre2013}). Con
\begin{equation}
  \omega = \frac{\sqrt{8 g h_0}}{a},
  \qquad
  A = \frac{a^2 - r_0^2}{a^2 + r_0^2},
\end{equation}
donde $r_0 = 0.8$\,m fija la amplitud de la oscilación, los campos son
\begin{align}
  h(x,y,t) &= h_0\left[
               \frac{\sqrt{1 - A^2}}{1 - A\cos(\omega t)} - 1
               - \frac{r^2}{a^2}\!
                 \left(\frac{1 - A^2}{(1 - A\cos(\omega t))^2} - 1\right)
              \right] - \zb(x,y),\\
  u(x,y,t) &= \frac{1}{2}\,
              \frac{\omega A \sin(\omega t)}{1 - A\cos(\omega t)}\,(x - x_c),\\
  v(x,y,t) &= \frac{1}{2}\,
              \frac{\omega A \sin(\omega t)}{1 - A\cos(\omega t)}\,(y - y_c),
\end{align}
con $h$ truncada a $0$ en la zona seca y $u = v = 0$ allí; la superficie libre
es $\eta = h + \zb$. La superficie permanece como un plano inclinado que
oscila en el tiempo.
